\documentclass[11pt,letterpaper]{article}

% --- packages ---------------------------------------------------------
\usepackage[margin=1in]{geometry}
\usepackage{amsmath,amssymb}
\usepackage{graphicx}
\usepackage{booktabs}
\usepackage{array}
\usepackage{hyperref}
\usepackage[dvipsnames]{xcolor}
\usepackage{fancyhdr}
\usepackage{titling}
\usepackage{enumitem}
\usepackage{mdframed}
\usepackage{textcomp}
\usepackage[normalem]{ulem}

\hypersetup{
  colorlinks=true,
  linkcolor=NavyBlue,
  citecolor=NavyBlue,
  urlcolor=NavyBlue
}

\pagestyle{fancy}
\fancyhf{}
\fancyhead[L]{\small Preprint, May 2026 --- v5}
\fancyhead[R]{\small Morin --- Multi-Beam Direct-Write EBL}
\fancyfoot[C]{\thepage}
\renewcommand{\headrulewidth}{0.3pt}

\setlength{\droptitle}{-2em}

\title{\Large \textbf{Multi-Beam Direct-Write Electron Lithography\\via Multi-Rate Electromagnetic Steering}\\[0.3em]
  \large v5 --- Open Software Stack and Final Integrated Release\\[0.2em]
  \normalsize \textit{Complete Open-Source Design Package: Simulated, Spec'd, Software-Defined}}
\author{Robert Gerald Morin\textsuperscript{1,*}\\[0.4em]
  {\small \textsuperscript{1}N\^{I}KI Nation Builders, Edmonton, Alberta, Canada}\\
  {\small \textsuperscript{*}Correspondence: \href{mailto:robmorin0@gmail.com}{robmorin0@gmail.com}}}
\date{Preprint v5 --- May 16, 2026\\[0.3em]
  {\small\textit{v3 (May 2026) published the cross-beam Coulomb analysis and 20\,\textmu m pitch reframe. v4 closed the principal engineering hurdles identified by external peer review of v3: Loeffler--Jansen stochastic blur, thermal management of $10^6$ deflection coils in UHV, magnetic crosstalk at 20\,\textmu m pitch, 400\,Tbps pattern-data streaming, resist outgassing, source-array uniformity, and active fiducial registration. v4.1 added independent Niki-run engineering sims confirming each subsystem claim and refining $\sigma_\text{total}$ from 8.0\,nm to 6.65\,nm. \textbf{v5 adds the open-source software stack:} a GDS-II $\to$ blanker-schedule pattern compiler with 3-stage hierarchical compression, a $10^6$-channel hard real-time control firmware closing the inner loop in 6\,\textmu s out of a 10\,\textmu s frame, and a tool-operator OS (PostgreSQL calibration database + Next.js GUI + SECS/GEM MES integration). With v5 the architecture is no longer just spec'd hardware: it is a complete, software-defined, end-to-end open-source design package. Each subsystem is detailed in companion specifications [16, S1--S11].}}}

\begin{document}
\maketitle
\thispagestyle{fancy}

% --- open-release banner ---------------------------------------------
\begin{mdframed}[linewidth=0.5pt,backgroundcolor=gray!4]
\textbf{Open Release.} This work is released under a Creative Commons Attribution 4.0 International License (CC BY 4.0). It may be freely copied, modified, built, commercialised, or extended by anyone, for any purpose, with attribution to the author. The author retains no patent claim on the architecture or its constituent ideas. By publication, all combinations and methods described herein enter the public domain as prior art, and may not be enclosed by subsequent intellectual-property filings by any party. The v5 software reference implementations are released under MIT (code) + CC BY 4.0 (spec).
\end{mdframed}

% --- abstract --------------------------------------------------------
\begin{center}\textbf{Abstract}\end{center}
\noindent
The v3 preprint (Morin 2026) proposed and quantitatively validated the cross-beam Coulomb perturbation envelope for an open, mask-free, multi-vendor multi-beam direct-write electron-beam lithography architecture targeting the mature-node semiconductor manufacturing segment. External peer review of v3 confirmed the cross-beam analysis but identified seven remaining engineering hurdles: (1) Loeffler--Jansen intra-beam stochastic blur; (2) thermal management of $10^6$ deflection coils dissipating $\sim$\,1\,kW in ultra-high vacuum without convective cooling; (3) magnetic crosstalk between coils at 20\,\textmu m pitch; (4) $\sim$\,400\,Tbps pattern-data streaming bandwidth; (5) resist outgassing under intense parallel exposure; (6) source-array tip-to-tip uniformity; (7) sub-nm placement registration across a $10^6$-channel array.

v4 closes each of these through subsystem-level engineering specifications. The architectural moves are:

\begin{itemize}[leftmargin=1.6em]
  \item \textbf{Cryogenic 77\,K column} with YBCO 2G coated-conductor deflection coils (zero $I^2R$ loss) and YBCO Meissner-effect magnetic shielding between coils ($>$\,80\,dB attenuation, eliminating crosstalk as a budget item).
  \item \textbf{Photonic data path} with 32 single-mode optical-fibre trunks $\times$ 96 wavelengths $\times$ 200\,Gbps PAM4 (608\,Tbps raw, 50\,Tbps after 10$\times$ hierarchical pattern compression), feeding cryo-CMOS DAC tiles inside the cold column.
  \item \textbf{Per-beam current de-rating} from v3's 5\,nA to v4's 2\,nA, driven jointly by Loeffler--Jansen budget (calibrated at 7\,nm via scaling from IMS Nano commercial data, and independently confirmed at 6.04\,nm RMS by Jansen-formula evaluation in v4.1) and wafer-side peak-flux thermal limits. This opens the validated node range to 50--180\,nm (with marginal access to 30\,nm) at $\sim$\,3$\times$ lower throughput than v3.
  \item \textbf{Three-tier active registration loop} with backscatter-electron fiducial probes ($10^3$ of $10^6$ beams used as registration sensors) and Zygo ZMI stage metrology, closing the placement budget at $\sigma_\text{total} \approx 6.65$\,nm (v4.1, refined from v4's 8.0\,nm via independent Boersch chromatic-blur evaluation).
  \item \textbf{AlN wafer chuck} with 5\,Torr backside He gas + Galden microchannel liquid loop + 36-zone TEC trim, holding $\pm 8$\,mK uniformity under 65\,W beam deposition.
  \item \textbf{Hierarchical UHV management}: cryosorption on the 77\,K column captures $\sim$\,80--90\,\% of resist outgassing on first impact, with conventional turbomolecular and ion pumping handling non-condensables.
\end{itemize}

\textbf{v5 closes the software loop.} The new \S 3.14 specifies three open-source subsystems: (i) a GDS-II $\to$ per-beam blanker-schedule \textbf{pattern compiler} with three-stage hierarchical compression (quadtree dedupe + RLE + Zstandard) achieving the v4 \S 3.6 10$\times$ compression target and emitting a self-delimited binary streaming format compatible with the photonic data path; (ii) a \textbf{$10^6$-channel real-time control firmware} stack (L0--L4) closing the inner deflection loop in 6\,\textmu s out of a 10\,\textmu s frame on per-tile cryo-CMOS RISC-V supervisors with Chisel-generated SystemVerilog; (iii) a \textbf{tool operating system} (calibration database, recipe/lot manager, Next.js operator GUI, SECS/GEM MES adapter, telemetry pipeline, alarm system) built on a pinned commodity stack (PostgreSQL 18 + MinIO + VictoriaMetrics + Tango Controls). Capital cost for the 10-tool-amortised software NRE: \$2--3\,M, raising the per-tool total to $\sim$\,\$40\,M.

The architecture is now \textbf{hardware-specified, sim-verified, and software-defined}, with every required library, framework, and tool pinned to a permissively-licensed open-source implementation. Updated capital cost estimate: \textbf{$\sim$\,\$40\,M per tool} (v5, +\$2.5\,M amortised software NRE on top of v4.1's \$37.6\,M), retaining a $\sim$\,10$\times$ cost advantage over EUV. Cycle-time at the v4 nominal 2\,nA derate: $\sim$\,8--17 wafers per hour through a 70-layer mature-node process, retaining order-of-magnitude advantages over single-beam EBL.

Every subsystem has $\geq 3$ independent commercial vendors across $\geq 2$ jurisdictions. The architecture is released openly as a \textbf{complete open-source design package}: paper + sim code + 11 subsystem specs + 3 software reference implementations. After this engineering pass, the remaining open work is process-development engineering (YBCO 20\,\textmu m patterning, cryo-CMOS DAC NRE, MEMS source array scale-up to $10^6$ tips) and software pre-validation on HIL (Stage A.5, \S 7). The fundamental physics and architecture are closed. \textbf{v5 is ready for Stage A prototype build, and the open-source package contains everything needed to begin.}

\medskip
\noindent\textbf{Keywords:} multi-beam EBL, cryogenic electron lithography, superconducting deflection coils, Meissner magnetic shielding, mature-node semiconductor manufacturing, sovereign semiconductor supply chain, silicon photonics, mask-free lithography, open-source EDA, GDS-II compiler, cryo-CMOS firmware

\medskip
\noindent\textbf{Classification:} Semiconductor manufacturing; Electron-beam lithography; Cryogenic engineering; Charged-particle optics; Silicon photonics; Real-time embedded systems; Open-source EDA

\bigskip\hrule\bigskip

% =====================================================================
\section{Carryforward from v3 and v4.1}

Sections 1, 2, 3.1--3.3, 6, 9 (CNT extensions), and 7 (geopolitical implications) of the v3 manuscript are carried into v5 without material change. The v3 cross-beam Coulomb analysis (\S 3.4) and the 20\,\textmu m nominal pitch derivation, along with the v4 subsystem engineering (\S 3.5--3.12) and the v4.1 Niki-run sim confirmations (\S 3.13 + Figures 5--8), are the foundation on which the v5 software stack is layered.

This v5 manuscript carries forward the full v4.1 hardware engineering content (\S 3.4.7 Loeffler--Jansen blur, \S 3.5 cryogenic column architecture, \S 3.6 photonic pattern data path, \S 3.7 per-beam blanker, \S 3.8 source array, \S 3.9 fiducial registration, \S 3.10 wafer thermal, \S 3.11 cold--warm UHV interface, \S 3.12 resist outgassing, \S 3.13 node-scale accessibility curve) and adds:

\begin{itemize}[leftmargin=1.6em]
  \item \textbf{\S 3.14} Software stack (new in v5): GDS-II pattern compiler, real-time control firmware, tool operating system
  \item \textbf{\S 4.3} updated capital cost breakdown including amortised software NRE
  \item \textbf{\S 7 Stage A.5} software pre-validation on HIL (new in v5)
  \item \textbf{\S 8} revised open questions (closing several v4 items, adding software-side items)
  \item \textbf{\S 9} updated conclusion
\end{itemize}

\S 5 (capital cost breakdown), \S 7 (validation pathway), \S 8 (remaining open questions), and \S 9 (conclusion) are revised.

% =====================================================================
\section{Loeffler--Jansen blur, calibrated from IMS Nano (\S 3.4.7)}

For a parallel electron beam at moderate current density, Jansen (1990) [10] derives the Loeffler trajectory-displacement blur in the ``weak interaction'' (Holtsmark) regime:

\begin{equation}
\sigma_{LJ} \;\propto\; \frac{I_b^{\,a}\,L^{\,b}}{V_0^{\,c}}, \qquad (a,b,c) \approx (2/3,\;3/2,\;3/2).
\end{equation}

The IMS Nanofabrication MBMW commercial mask-writer (262,000 beams, 50\,kV, $\sim$\,50\,mm drift, 0.5\,nA/beam in production mode) reports a total placement accuracy of $<$\,2\,nm [5, 17]. Attributing approximately 1\,nm of that budget to intra-beam Coulomb effects (the remainder being stage metrology, optical aberration, source ripple, and registration), and scaling to v4 nominal:

\begin{equation}
\sigma_{LJ}^{\,v4}\bigl(I_b = 2\,\text{nA},\;L=100\,\text{mm}\bigr) \approx 1\,\text{nm}\times\bigl(2/0.5\bigr)^{2/3}\times\bigl(100/50\bigr)^{3/2} \approx 1\,\text{nm}\times 2.52\times 2.83 \approx 7.1\,\text{nm}.
\end{equation}

With the uncertainty band on the IMS attribution and the scaling exponent ($a \in [0.62, 0.70]$), the prediction range is \textbf{$\sigma_{LJ} \in [5, 10]$\,nm at $I_b = 2$\,nA}. This is inside the conservative literature band cited by external review (5--20\,nm) and supports the v3 \S 3.4.6 open-budget treatment now closed at the v4 operating point.

\textbf{v4.1 update}: the IMS-Nano-anchored estimate (7\,nm) and independent Jansen-formula evaluation (6.04\,nm RMS at v4 nominal, see Figure~\ref{fig:phase2}) converge to within 15\%. The Jansen closed-form evaluation uses $\sigma_{LJ}[\text{nm}] = K \cdot I_b[\text{nA}]^{2/3} \cdot L[\text{mm}]^{3/2} \cdot V_0[\text{kV}]^{-3/2}$ with $K = 1.587$ fitted against the IMS Nano anchor (0.5\,nA / 50\,mm / 50\,kV $\to$ 1\,nm attributed), and the uncertainty band ($K \in [1.2, 2.0]$) gives an RMS range 4.57--11.41\,nm. Two independent methods now bracket the same value.

The current de-rating from v3's 5\,nA to v4's 2\,nA serves two roles simultaneously: it brings $\sigma_{LJ}$ from the predicted 13\,nm at 5\,nA down to 7\,nm at 2\,nA (LJ scales as $I_b^{2/3}$), and it brings peak local wafer flux below the resist-melt threshold (see \S 3.10).

Full derivation, scaling tables, and calibration assumptions are in [S1, \texttt{v4\_loeffler\_jansen.md}].

\begin{figure}[t]
\centering
\includegraphics[width=0.95\textwidth]{figure_phase2.pdf}
\caption{\textbf{v4.1 Phase 2 (Niki-run) independent Jansen-formula evaluation of Loeffler--Jansen blur, and refined placement-error budget.} Left: closed-form Jansen evaluation calibrated to the IMS Nano anchor returns $\sigma_{LJ} = 6.04$\,nm RMS at v4 nominal (2\,nA / 100\,mm / 50\,kV), converging within 15\% of the IMS-scaling estimate of 7\,nm. Uncertainty band $K \in [1.2, 2.0]$ gives 4.57--11.41\,nm RMS. Right: refined placement-error quadrature budget at v4.1, with Boersch chromatic blur revised from a conservative 3\,nm (v4) to 1.35\,nm (realistic $C_c = 50$\,\textmu m), dropping $\sigma_\text{total}$ from 8.0\,nm to 6.65\,nm.}
\label{fig:phase2}
\end{figure}

% =====================================================================
\clearpage
\section{Cryogenic column architecture (\S 3.5)}

The v3 architecture left the column at ambient temperature. At $N=10^6$ normal-metal deflection coils driven at the 1\,mT envelope of \S 3.1 and switched at 10--100\,kHz, resistive dissipation is $\sim$\,1\,kW concentrated in a 20\,mm $\times$ 20\,mm $\times$ 100\,mm column volume. Without convective cooling (UHV constraint), the column would run away thermally within milliseconds. The v3 manuscript flagged ``mechanical and thermal stability'' as an open question but did not solve it; the v3 peer review elevated it to a tool-killer.

The v4 architectural response is to operate the column at 77\,K with high-temperature superconducting (HTS) coils:

\begin{itemize}[leftmargin=1.6em]
  \item \textbf{Coil material}: YBCO 2G coated conductor on \textbf{silicon substrate} (revised from sapphire after CTE reconciliation with the silicon-substrate blanker die, \S 3.7 --- sapphire-vs-silicon CTE mismatch would have given $\sim$\,200\,\textmu m differential contraction over the 20\,mm column on 295\,K $\to$ 77\,K cooldown), planar Helmholtz-pair lithographed spirals at 20\,\textmu m pitch. Critical current density at 77\,K, self-field, is $J_c \approx 2 \times 10^{10}$\,A/m$^2$ (conservative), supporting the 1\,mT v3 envelope with $\sim$\,48$\times$ current margin. NbTi (9.2\,K), Nb$_3$Sn (18\,K), and MgB$_2$ (39\,K) are all excluded by the 77\,K operating point; only YBCO and similar HTS materials work.

  \item \textbf{Magnetic crosstalk solution}: a 10\,\textmu m-thick YBCO Meissner screen between adjacent coils expels magnetic flux completely (Meissner effect). Theoretical attenuation is $e^{-50} \approx 440$\,dB; defect-limited real-world performance from SQUID and dark-matter experimental literature is 80--120\,dB. Even at the pessimistic 80\,dB, neighbour crosstalk at 20\,\textmu m spacing is suppressed from a v3 worst-case $-30$\,dB (geometry alone) to $-110$\,dB --- below the 16-bit DAC LSB at $-96$\,dB. \textbf{Crosstalk is removed from the placement-error budget.} v4.1 qualifier: verified achievable with research-grade PLD YBCO (defect density $\geq 10^{10}$) giving 100\,dB attenuation comfortably below the 16-bit DAC LSB (96\,dB); commercial 2G coated conductor gives a marginal 80\,dB. See Figure~\ref{fig:meissner}.

  \item \textbf{Cryocooler (v4.1 revised)}: the integrated 77\,K cold-mass dissipation is $\sim$\,2\,W (coil cold zone) plus $\sim$\,150\,W (electronics cold zone, dominated by 10$^6$-channel cryo-CMOS DAC array at 100\,\textmu W each = 100\,W, plus 32-trunk photonic receivers at 1.5\,W each = 48\,W from \S 3.6), for an integrated load of \textbf{$\sim$\,151\,W} (with 50\% engineering margin: 227\,W). The corrected v4.1 cryocooler is a \textbf{dual-zone configuration}: Cryomech PT-180 pulse-tube (50\,W lift at 77\,K, $\sim$\,\$45\,k) on the coil cold zone, and Cryomech AL-200 (200\,W lift class at 77\,K, $\sim$\,\$80\,k) on the electronics cold zone, with N+1 redundancy giving $\sim$\,\$200\,k total cryocooler subsystem. Capital impact: \textbf{+\$170\,k = +0.5\%}, raising the cryocooler line from \$1.5\,M to \$1.67\,M and the total from \$37.4\,M to \$37.6\,M. Pulse-tube cycle is selected over Gifford-McMahon to reduce vibration coupling to the column (sub-100\,nm GM displacement vs sub-10\,nm pulse-tube). See Figure~\ref{fig:thermal}.

  \item \textbf{Cryo-CMOS DACs}: the per-coil 16-bit current DACs are realised in GlobalFoundries 22FDX SOI process at 77\,K, where threshold-voltage scatter is reduced and switching speed is $\sim$\,2$\times$ higher than at 295\,K. Per-DAC dissipation $\leq 100$\,\textmu W; total DAC array load at 77\,K is 100\,W, handled by the AL-200 electronics-zone cryocooler.
\end{itemize}

Full derivation, vendor list, AC-loss budget, and quench-protection architecture are in [S2, \texttt{v4\_cryocolumn.md}]. The cryocooler is a commercial commodity (multi-vendor: Sumitomo, Cryomech, Brooks); YBCO 2G tape is available from SuperPower, American Superconductor, SuperOx, Furukawa. No new single-vendor chokepoint.

\begin{figure}[t]
\centering
\includegraphics[width=0.95\textwidth]{figure_phase3_thermal.pdf}
\caption{\textbf{v4.1 cryogenic column integrated thermal budget at 77\,K (Niki-run Phase 3 sim).} Stacked contributions sum to $\sim$\,151\,W; with 50\% engineering margin the design point is 227\,W. v4.1 cryocooler stack: PT-180 (50\,W @ 77\,K) on coil cold zone + AL-200 (200\,W @ 77\,K) on electronics cold zone, both N+1 redundant.}
\label{fig:thermal}
\end{figure}

\begin{figure}[t]
\centering
\includegraphics[width=0.95\textwidth]{figure_phase3_meissner.pdf}
\caption{\textbf{v4.1 YBCO Meissner shield attenuation vs film quality (Niki-run Phase 3 sim).} Theoretical maximum 434\,dB; geometric edge leakage 20.8\,dB; defect-limited residual set by film quality. Commercial 2G coated conductor gives 80\,dB (marginal); research-grade PLD YBCO gives 100\,dB (comfortable); SQUID-grade gives 120\,dB (24\,dB margin). v4.1 specifies research-grade-or-better PLD YBCO.}
\label{fig:meissner}
\end{figure}

% =====================================================================
\section{Photonic pattern data path (\S 3.6)}

The v3 manuscript identified ``mask-free pattern data handling'' as an open question. The naive uncompressed pixel-rate equivalent for the v4 architecture is $\sim$\,400\,Tbps; the v3 peer review elevated this to a feasibility question.

The v4 data path topology is three-tiered:

\begin{itemize}[leftmargin=1.6em]
  \item \textbf{L0 sync}: 1\,Gbps clock and frame synchronisation
  \item \textbf{L1 deflection}: $10^6 \times 16$\,bits $\times$ 2 axes $\times$ 10\,kHz = 320\,Gbps to per-coil DACs
  \item \textbf{L2 blanker + pattern}: $\sim$\,50\,Tbps to per-beam blanker layer, after 10$\times$ hierarchical pattern compression (IMS-Nano-style region encoding + LZMA residuals)
\end{itemize}

Aggregate after compression: \textbf{$\sim$\,50\,Tbps sustained}, carried over \textbf{32 single-mode fibre trunks}, each running 96 wavelengths $\times$ 200\,Gbps PAM4 = 19.2\,Tbps per trunk. Raw fibre capacity 614\,Tbps; loaded at 12\,\% gives substantial headroom.

In-column reception: UHV-compatible fibre feedthroughs (Accu-Glass, Thorlabs, Larson; multi-vendor); cryo-compatible InP photodiode arrays at 77\,K with improved responsivity and dark current; per-coil decode/buffer ASICs in cryo-CMOS 22FDX.

Total in-column electronic dissipation (receiver + DAC + buffer + decode): 150\,W at 77\,K, fitted to the AL-200 electronics-zone cryocooler stage (\S 3.5). Recurring per-tool cost \$4.5\,M; first-article including ASIC NRE \$12.5\,M; 10-tool amortised \$5.3\,M.

Full topology, fibre/transceiver vendor list, compression engine, and ASIC dual-source strategy in [S3, \texttt{v4\_datapath.md}].

% =====================================================================
\section{Per-beam blanker subsystem (\S 3.7)}

The middle loop in the multi-rate control hierarchy (v3 \S 2.3) blanks beams ON/OFF in synchrony with pattern data. At $10^6$ beams $\times$ 1\,kHz update $\times$ 1\,bit = 1\,Gbps base rate (with optional 4-bit dose modulation at 16$\times$ rate for greyscale).

The v4 blanker is an electrostatic plate array fabricated as a three-layer hybrid-bonded CMOS-MEMS stack:

\begin{itemize}[leftmargin=1.6em]
  \item Each beam passes between a pair of plates, \textbf{50\,\textmu m long with 10\,\textmu m gap}, requiring \textbf{40\,V} switched voltage to produce a 2\,mrad deflection --- enough to clear the stop aperture below. RC bandwidth $\sim$\,1\,ns (1\,GHz), 6 decades margin over the 1\,kHz blanker spec.
  \item Three-layer hybrid-bonded stack: (i) plate die (silicon, fabricated with deep RIE etched plates), (ii) HV-CMOS driver die (silicon, per-blanker level-shifter + switch), (iii) Mo+graphite stop aperture (separate water-cooled assembly, \textit{outside} the cold mass).
  \item Column footprint at 20\,\textmu m pitch over $10^6$ blankers: 20 $\times$ 20\,mm --- single die fits within the standard $\sim$\,26 $\times$ 33\,mm reticle limit; \textbf{no tiling required}.
  \item Cold-side power: 5\,W dynamic switching + 10\,W serial-link receivers + 7\,W margin = \textbf{22\,W at 77\,K}.
  \item Subsystem cost: \$185\,k production + \$4.3\,M NRE $\to$ \$0.19 per blanker amortised over $10^6$ units. Under 0.5\,\% of total tool capital.
\end{itemize}

Full details, sapphire-vs-silicon CTE reconciliation, charging mitigation, and 4-bit greyscale dose option in [S4, \texttt{v4\_blanker.md}].

% =====================================================================
\section{Source array and per-beam current control (\S 3.8)}

The v3 manuscript flagged ``source array uniformity'' as an open question: field-emitter tip-to-tip current variations of 10--30\,\% require per-beam closed-loop control.

The v4 source is a \textbf{MEMS CNT/Spindt cold-field-emitter array} at 20\,\textmu m pitch. ZrO/W Schottky emitters are excluded by the 20\,\textmu m pitch constraint (each Schottky emitter assembly is 3+ orders larger). Vendor architecture: NSC for the tip IP, CEA-Leti or imec for the foundry.

Per-tip current sensing: a thin annular sense electrode at each aperture taps $\sim$\,1\,\% of the beam current --- 50\,pA at 5\,nA, 20\,pA at the v4 2\,nA derated current. Cryo-CMOS transimpedance amplifier (TIA) at 77\,K with 5\,fA/$\sqrt{\text{Hz}}$ noise floor (Caeleste, Teledyne). 12-bit extractor DAC per tip, 50\,mV resolution against $\sim$\,1\,V Fowler-Nordheim sensitivity, 1\,kHz loop bandwidth.

Per-channel electronics cost: $\sim$\,\$0.65 (12-bit DAC \$0.05, HV level-shifter \$0.20, share of readout ASIC + FPGA \$0.30, share of digital backplane \$0.10). Total for $10^6$ channels: \textbf{\$650\,k}. The architecture's premise that ``per-beam everything'' is now economic is validated.

Full architecture in [S5, \texttt{v4\_source\_and\_registration.md}] Part A.

% =====================================================================
\section{Fiducial registration subsystem (\S 3.9)}

The v3 manuscript flagged ``mechanical and thermal stability'' as an open question. Across a 20\,mm column array, even sub-\textmu m thermal drift creates registration errors that destroy the 6\,nm placement budget.

The v4 active-registration loop is \textbf{three-tiered}, matching the multi-rate control philosophy:

\begin{itemize}[leftmargin=1.6em]
  \item \textbf{Inner (10\,kHz, per-beam)}: backscatter-electron (BSE) detection on etched fiducial marks on each die. $10^3$ of the $10^6$ beams (0.1\,\%) are used as registration probes.
  \item \textbf{Middle (100\,Hz, per-column-subarray)}: aggregate per-row drift, fed back as a column-row trim DAC.
  \item \textbf{Outer (1\,Hz, column-wide)}: stage interferometric metrology (Zygo ZMI differential interferometer) provides absolute-frame anchoring at 0.1\,nm/Hz resolution.
\end{itemize}

Feedback to the deflection coils via a 12-bit $\pm 500$\,nm trim DAC summed analog into the existing 16-bit per-beam deflection DAC. Vendor list: El-Mul, Hamamatsu (BSE+SiPM); Zygo, Heidenhain (interferometry); Instrumentation Technologies Libera (BPM readout). Subsystem cost \$15\,M total, dominated by the high-precision interferometric stage metrology. Full architecture in [S5, \texttt{v4\_source\_and\_registration.md}] Part B.

% =====================================================================
\section{Wafer-side thermal management (\S 3.10)}

A $10^6$-beam array at v3's 5\,nA / 50\,kV deposits $\sim$\,250\,W of beam power onto the wafer when all beams are unblanked. Averaged over the 26\,\% write-cycle duty, this is 65\,W steady-state; the \textbf{peak local flux} beneath a 20\,mm column footprint is $6.25 \times 10^5$\,W/m$^2$, producing a transient $\Delta T$ of $\sim$\,96\,K in 0.26\,s. \textbf{This melts the resist.}

The v4 architecture selects \textbf{per-beam current derating to 2\,nA} as nominal. At 2\,nA / $10^6$ beams / 20\,mm column, peak local flux is $2.5 \times 10^5$\,W/m$^2$ --- manageable with active chuck cooling.

The wafer chuck architecture:

\begin{itemize}[leftmargin=1.6em]
  \item \textbf{AlN ceramic chuck body} (high thermal conductivity, low CTE, UHV-compatible).
  \item \textbf{5\,Torr backside He gas} in a 20\,\textmu m gap between chuck and wafer underside (standard semiconductor-lithography practice).
  \item \textbf{Galden HT-135 perfluoropolyether liquid loop} in vacuum-brazed microchannels internal to the chuck.
  \item \textbf{36-zone Peltier (TEC) trim} under the chuck for fine spatial uniformity ($\pm 8$\,mK across the 300\,mm wafer).
\end{itemize}

Subsystem cost \$450\,k. Wafer-to-wafer cycle (thermal-soak limited): 6.7\,min/wafer single-loadlock, 1\,min/wafer dual-loadlock with pre-heat. Tool throughput at the v4 nominal: \textbf{$\sim$\,8--17\,wph through a 70-layer mature-node process}. Full architecture, equations, and vendor list in [S6, \texttt{v4\_wafer\_thermal.md}].

% =====================================================================
\section{Cold--warm UHV interface (\S 3.11)}

Between the 77\,K column and the 295\,K wafer chuck is a 50\,mm UHV gap that electrons must traverse without deflection. The radiative load from the warm wafer onto the cold column must fit within the column thermal budget; mu-metal magnetic shielding must compensate Earth's field and any stray laboratory magnetic field; and mechanical decoupling must prevent stage vibrations from coupling into the column.

The v4 interface includes a view factor of 0.560 for two coaxial discs giving 1.5\,mW direct radiative load with polished-gold liner; an intermediate 150\,K radiation shield with 10-layer aluminised-Kapton MLI; three-layer mu-metal magnetic shielding with Cryoperm 10 inner can (Vacuumschmelze) giving $\sim$\,85\,dB DC attenuation; 1100\,kg granite base + 4$\times$ Newport S-2000A active pneumatic isolators; and differential pumping with 100\,mm CF turbo port (Pfeiffer HiPace 700, N+1 redundant) for 600$\times$ pressure step.

Total interface cost: \textbf{\$1.22\,M}. Full architecture in [S7, \texttt{v4\_cold\_warm\_interface.md}].

% =====================================================================
\section{Resist outgassing management (\S 3.12)}

The v3 peer review flagged that high-density beam exposure of chemically-amplified resists (CARs) generates volatile species that contaminate the cold column. v4 resolves this through resist selection and pumping topology:

\begin{itemize}[leftmargin=1.6em]
  \item \textbf{Resist family}: \textbf{HSQ (Dow XR-1541)} for prototyping; \textbf{Inpria Sn-cluster metal-oxide resist} for production (50--100$\times$ lower outgassing than CAR).
  \item \textbf{Outgassing burst} during 0.26\,s write: $\sim$\,7.1 $\times 10^{17}$ molecules per 300\,mm wafer.
  \item \textbf{Cryosorption on the 77\,K column annulus} pumps $\sim$\,2000\,L/s for condensable organics with sticking probability $s \approx 0.5$, absorbing 80--90\,\% of the outgassing burst on first impact.
  \item \textbf{Conventional pumping}: 1$\times$ Pfeiffer HiPace 2300 magnetic-levitation turbomolecular pump + 2$\times$ Gamma 500\,L/s ion pumps handle the non-condensable fraction.
  \item \textbf{Regeneration cycle}: monthly warm-up of cryosorbed surfaces releases $\sim$\,18\,g/year accumulated organic ice; 2.25\,h/month downtime = \textbf{0.3\,\%}, under the 1\,\% budget.
\end{itemize}

Pumping subsystem cost \$136\,k. Full vendor list and contamination-buildup model in [S8, \texttt{v4\_resist\_outgassing.md}].

% =====================================================================
\section{Node-scale accessibility curve (\S 3.13, new in v4.1)}

The v4.1 Niki-run Phase 3 sim plots $\sigma_\text{total}$ as a function of per-beam current $I_b$ with achievable lithographic node overlaid (Figure~\ref{fig:nodescale}). The Loeffler--Jansen and Boersch contributions both scale with $I_b$; in the $I_b \to 0$ limit they vanish and the only remaining contributions are the architectural floors:

\begin{equation}
\sigma_\text{total}(I_b \to 0) \;\to\; \sqrt{\sigma_\text{stage}^2 + \sigma_\text{column}^2} \;=\; \sqrt{1^2 + 1^2} \;=\; \mathbf{1.41\,\text{nm}}.
\end{equation}

This is the architecture's \textbf{hard floor}: $\sigma_\text{total}$ cannot drop below $\sim 1.41$\,nm regardless of current de-rating, which translates to a \textbf{$\sim 7$\,nm node accessibility limit}.

The current/node tradespace at v4.1 operating points:

\begin{center}
\begin{tabular}{rrrr}
\toprule
$I_b$ (nA) & $\sigma_\text{total}$ (nm) & Min accessible node & Throughput vs v4 \\
\midrule
5.0 (v3 nominal) & 12.5 & 63\,nm & 2.5$\times$ \\
\textbf{2.0 (v4 nominal)} & \textbf{6.66} & \textbf{33\,nm} & \textbf{1.0$\times$} \\
1.0 & 4.27 & 21\,nm & 0.5$\times$ \\
0.5 & 2.88 & 14\,nm & 0.25$\times$ \\
0.1 & 1.65 & 8\,nm & 0.05$\times$ \\
$\to 0$ & 1.41 (floor) & 7\,nm & $\to 0$ \\
\bottomrule
\end{tabular}
\end{center}

Headline implications: v4.1 architecture is \textbf{physics-capable of single-digit-nm nodes} at very low current, and \textbf{production-capable at 30--180\,nm} at 2\,nA nominal. The current/node tradeoff is \textbf{continuous} --- no sharp cliff --- which lets a single tool serve multiple market segments by changing only the per-beam current setting.

\begin{figure}[t]
\centering
\includegraphics[width=0.95\textwidth]{figure_phase3_node_scale.pdf}
\caption{\textbf{v4.1 node-scale accessibility curve (Niki-run Phase 3 sim).} $\sigma_\text{total}$ as a function of per-beam current $I_b$ with the architectural floor $\sigma_\text{floor} = 1.41$\,nm shown as the asymptote at $I_b \to 0$. The current/node tradeoff is continuous, allowing a single tool to serve multiple node markets by changing only $I_b$.}
\label{fig:nodescale}
\end{figure}

% =====================================================================
\clearpage
\section{Software stack (\S 3.14, new in v5)}

The v3 and v4 manuscripts both flagged software as a partially-specified item: v3 \S 3.4.9 noted that ``pattern data delivery and the offline compile chain are routine engineering on the standard EBL playbook,'' and v4 \S 3.6 closed the photonic transport but left a hand-off note in \texttt{v4\_datapath.md} \S 7 that ``\S 9 v4\_software.md: GDS-II $\to$ segment-stream compiler is the long pole for tool turn-around.'' The v5 software stack is the response.

Three open-source subsystems, each fully specified to a runnable code skeleton, with vendor / library pinning matching the niki-os monorepo conventions:

\begin{enumerate}[leftmargin=1.6em]
  \item \textbf{Pattern compiler} (\S 3.14.1, [S9] \texttt{v5\_software\_compiler.md}). Offline: GDS-II $\to$ per-beam blanker schedule + per-coil deflection trajectory + binary streaming format. Target: 10$\times$ hierarchical compression to match the 50\,Tbps loaded data-path budget.
  \item \textbf{Real-time control firmware} (\S 3.14.2, [S10] \texttt{v5\_software\_firmware.md}). Hard real-time: $10^6$-channel L0--L4 stack on per-tile cryo-CMOS RISC-V supervisors. Target: close inner deflection loop in 6\,\textmu s out of a 10\,\textmu s frame.
  \item \textbf{Tool operating system} (\S 3.14.3, [S11] \texttt{v5\_software\_operating\_system.md}). Soft real-time / supervisory: calibration database, recipe/lot manager, operator GUI, SECS/GEM MES adapter, telemetry pipeline, alarm system. Target: standard fab MES integration, multi-tenant customer API, full audit trail.
\end{enumerate}

All three are released under MIT (code) + CC BY 4.0 (spec) with no patent claim. Reference implementations include buildable code skeletons (Python + Rust + SystemVerilog + Chisel + cocotb + Next.js + SQL) sufficient that a competent team can extend them to production without re-deriving architecture decisions.

\subsection{Pattern compiler (\S 3.14.1)}

The compiler sits between two industry-defined interfaces. \textbf{Upstream}: the customer's tape-out database --- GDS-II (Calma 1980s binary, de-facto standard) or its modern successor OASIS (SEMI P39, $\sim$10$\times$ smaller for the same layout). \textbf{Downstream}: the per-beam segment stream defined by the v4 photonic data path (\S 3.6), namely a sequence of tile-keyed binary blocks consumed by the in-column pattern-segment expander ASIC (64\,kB FIFO per beam, RLE + hierarchical-fill decoder).

\textbf{Key design choices.}

\begin{itemize}[leftmargin=1.6em]
  \item \textbf{Offline compile, real-time stream.} A 70-layer mask set takes hours to compile and is then re-streamed at run time from a 1--2\,PB SSD pool. Re-streaming is real-time (50\,Tbps sustained); compilation is not. This is the same separation MAPPER and IMS Nano use in their production flows.
  \item \textbf{Library selection.} \texttt{gdstk} (BSL-1.0, C++ core with Python bindings, $\sim$10$\times$ faster than \texttt{gdspy}) for the production ingest path; \texttt{klayout} (GPL-3) used out-of-tree as a DRC oracle but kept out of the MIT-licensed main distribution. OASIS supported via the same \texttt{gdstk} library.
  \item \textbf{Three-stage compression cascade.} (1) quadtree region-dedupe with content-addressed leaves (xxHash or BLAKE3), 3--5$\times$ on logic / 20--50$\times$ on memory; (2) byte-aligned RLE within leaves, 1.5--2$\times$; (3) Zstandard entropy coding (BSD-licensed, FPGA-accelerated via AMD/Xilinx Versal + Eideticom NoLoad), 1.5--2$\times$. Cascade product 7--20$\times$ on logic, 45--200$\times$ on memory. Meets the v4 \S 3.6 10$\times$ target with margin.
  \item \textbf{Hot-reload calibration.} Per-beam calibration corrections (offset, gain, distortion polynomial coefficients) are stored in a 96\,MB table that the compiler injects directly into the deflection trajectory and dose map during rasterisation. At runtime the table is updated at 1\,kHz from the in-column metrology loop (BSE fiducials, \S 3.9); the streamer holds two tables and atomically swaps on a frame boundary. The compiled pattern is \textbf{not} recompiled --- calibration is late-binding.
  \item \textbf{Binary streaming format.} Self-delimited macro-tile records with 64-byte header (\texttt{MBMW} magic, version, tile-id, dimensions, dose bits, compression flag, BLAKE3 content digest) + leaf dictionary + Zstandard-compressed payload. Intentionally similar to IMS Nano's published shot-data container format so commercial mask shops can adopt it without retraining.
  \item \textbf{Toolchain pinning.} Python 3.12 + Rust 1.80 via PyO3 (hot loops); \texttt{uv} + \texttt{pyproject.toml} build; \texttt{pytest} + \texttt{hypothesis} tests; \texttt{ruff} + \texttt{pyright --strict} lint; \texttt{mkdocs-material} docs. Matches niki-os monorepo conventions.
\end{itemize}

\textbf{Production scale.} A single 32-core / 256\,GB / 4\,TB NVMe workstation compiles one mask layer in 10--60 minutes; a 4-node cluster compiles a 70-layer mask set in $<$\,8\,h overnight. The 32-server streamer pool (sized for re-streaming, not compilation) sustains the 50\,Tbps column-boundary rate during exposure with a $<$\,500\,\textmu s end-to-end tile-to-blanker-plate latency.

\textbf{Code skeleton.} Reference repo \texttt{multibeam-ebl-compiler/} contains a $\sim$90-line Python skeleton covering GDS-II ingest, hierarchy flatten, and 30\,nm-grid rasterisation of a single tile, sufficient to round-trip a synthetic GDS test case in $<$\,10 minutes from clone. Production rasteriser hot loop is in Rust via PyO3. Full details, repo layout, format spec, and 7 open questions for v5.1+ in [S9] \texttt{v5\_software\_compiler.md}.

\subsection{Real-time control firmware (\S 3.14.2)}

The firmware closes the three nested loops introduced in v3 \S 2.3 (inner Lorentz deflection at 10--100\,kHz, middle blanker + dose at 1\,kHz, outer wafer stage at 1--10\,Hz) plus the registration trim (100\,Hz nominal, summed into the inner loop) across $10^6$ channels. The architecture is layered L0--L4 to align with the v4 \S 3.6 photonic tiers.

\textbf{Key design choices.}

\begin{itemize}[leftmargin=1.6em]
  \item \textbf{L0--L2 in RTL on cryo-CMOS, L3--L4 in software.} The deterministic data path (photonic PHY $\to$ MAC + FEC + framing $\to$ per-channel pattern decode $\to$ per-coil FIFO) lives entirely in RTL on the 1000 cryo-CMOS tile ASICs. L3 (per-tile supervisor, RV32IMC Ibex RISC-V core running Zephyr RTOS) handles fault recovery, BIST, calibration-time configuration. L4 (warm-side x86 + Linux + Tango Controls) handles tool start/stop, recipe load, Kalman fusion for registration. The inner-loop data path \textbf{never} traverses the RISC-V core.
  \item \textbf{Per-coil inner-loop FSM.} 1024 instances per tile, generated from a single parametric Chisel module. Five-state machine (IDLE $\to$ ARM $\to$ ACQ $\to$ ACT $\to$ SETTLE) with PI control (Kp = 0.25, Ki = 0.05 in Q24 fixed-point, bounded integrator to prevent windup), backscatter sense-electrode current monitor, and a one-way FAULT $\to$ BLANK\_PERMANENT trap for $>$\,5\,\% sense-vs-target deviation over 3 consecutive frames.
  \item \textbf{Latency budget.} Photonic L0/L1 receive $<$\,100\,ns; L1 MAC parse + FEC $<$\,200\,ns; L2 decode + dispatch $<$\,500\,ns; L3 PI compute $<$\,200\,ns; DAC slew $<$\,5\,\textmu s. Total \textbf{$<$\,6\,\textmu s out of the 10\,\textmu s frame, 0.6 utilisation, 4\,\textmu s margin} for fault handling and trim updates. The DAC slew is the dominant budget item and drives the cryo-CMOS DAC architecture choice (slew-rate-enhanced segmented architecture).
  \item \textbf{Synchronisation.} A single femtosecond-stable master clock (Menlo Systems FC-1500 or equivalent OFC) is encoded onto a dedicated L0 photonic lane and broadcast to every tile, identical in topology to JESD204C deterministic latency. The ZMI stage-position counter is continuously broadcast at 1\,MHz on a second L0 lane (32-bit sub-nm units). Every blanker event carries a 64-bit wafer-time stamp.
  \item \textbf{Registration Kalman fusion.} 6-state Kalman filter per column subarray ($\sim$\,$10^4$ beams per subarray, 100 subarrays total) on the warm-side L4 supervisor. States $\{\Delta x, \Delta y, \Delta\theta, \Delta\text{scale}, \Delta\text{skew}, \Delta\text{focus}\}$. Measurement updates from BSE probes at 10\,kHz, ZMI at 1\,kHz, thermal model at 1\,Hz. Fused state propagated to per-coil trim DAC offsets at 100\,Hz.
  \item \textbf{Hot redundancy.} Hot-spare beams allocated at column-design time (typically 1--5\,\%); faulted beams permanently blanked and deflected to graphite stop apertures. The L4 supervisor reallocates pattern data away from dead channels to designated hot-spare beams. Tile-level failure: 1024 beams ($\sim$\,0.1\,\% of array) routed to the spare pool, tool continues at degraded throughput until the next maintenance window.
  \item \textbf{Toolchain pinning.} Chisel 6 (MIT, Scala) for parametric HDL, with Amaranth (Python) as fallback; FIRRTL $\to$ SystemVerilog emission; Verilator + cocotb (LGPL + BSD) simulation; Verible lint; SymbiYosys formal; Vivado 2024.x for FPGA validation builds on VCU128 (Versal Premium); Cadence Innovus or Synopsys ICC2 for ASIC P\&R (proprietary, under foundry MPW contract); GCC + LLVM riscv-gnu-toolchain; Zephyr 4.x RTOS (Apache 2, primary; FreeRTOS as fallback); Tango Controls (LGPL) warm-side; gRPC + protobuf for the 10\,GbE L4 $\leftrightarrow$ L3 control plane; Nix flake for reproducible builds.
\end{itemize}

\textbf{ASIC impact.} The per-tile cryo-CMOS ASIC adds a $\sim$50\,k-gate RV32IMC core + 320\,kB SRAM to the existing 1024-channel DAC tile. This fits inside the existing $\sim$2\,mm$^2$ per-tile area budget from v4 \S 3.6; the \$200/tile production cost and \$8\,M NRE (dual-sourced via Imec + Fraunhofer IPMS) are unchanged by the firmware additions.

\textbf{Code skeleton.} Reference repo \texttt{multibeam-ebl-firmware/} contains a $\sim$45-line cocotb testbench (Python) and a $\sim$45-line SystemVerilog inner-loop FSM module, each runnable under Verilator inside the supplied Nix flake. Full repo, MIT + LGPL split, vendor wrappers, and 8 open questions for v5.1+ in [S10] \texttt{v5\_software\_firmware.md}.

\subsection{Tool operating system (\S 3.14.3)}

The tool OS is the soft-real-time supervisory layer above the firmware: calibration archive, recipe/lot management, operator GUI, telemetry pipeline, alarm system, and external integration (MES, customer-facing REST API). It runs on commodity Linux (Ubuntu 24.04 LTS) on x86\_64-v3 control-room servers --- typically 2$\times$ Dell R760 or equivalent (dual Xeon Gold, 256\,GB ECC, 8$\times$ 3.84\,TB NVMe in RAID-10).

\textbf{Key design choices.}

\begin{itemize}[leftmargin=1.6em]
  \item \textbf{Five-plane architecture.} Field plane (cryo-CMOS firmware, 1\,\textmu s--100\,\textmu s, see \S 3.14.2); cell controller (per-tile RTOS, 100\,\textmu s--10\,ms); \textbf{supervisory} (Linux + Rust services, 10\,ms--1\,s, this subsystem); \textbf{operator/process} (Next.js 15 SPA, 100\,ms-human); \textbf{enterprise} (Postgres + VictoriaMetrics + MinIO + Grafana, seconds-days). The supervisory plane is the seam between firmware telemetry (Cap'n Proto over 10\,GbE TSN, $\sim$\,30\,MB/s aggregate) and operator/MES intent (REST + WebSocket).
  \item \textbf{Calibration database.} PostgreSQL 18 as the index of record (metadata under 1\,TB even at multi-year archive); MinIO (S3-compatible, self-hosted) for the per-beam coefficient blobs ($\sim$\,100\,MB/snapshot, replicated to a secondary site); VictoriaMetrics 1.110 cluster for the 100\,GB/day time-series telemetry firehose (Apache 2 license, single-binary, PromQL-compatible). A snapshot is the atomic unit: it pins one row in every coefficient table to a single UUID v7, signed by the process engineer, with an instant-rollback pointer table. Calibration procedures (deflection-coil offset, blanker delay, tip emission profile, thermal map, registration zero, source uniformity, full PM) run as versioned recipes with acceptance bands.
  \item \textbf{Recipe + lot model.} Recipes are directories in a Gitea-hosted Git repo, content-addressed by commit SHA --- production lots always reference an immutable hash. Edits go through pull request with mandatory second-engineer review. Lot rows in Postgres reference the recipe SHA and the active calibration snapshot UUID; per-wafer rows hold start/end timestamps, yield, and MinIO telemetry-blob pointers.
  \item \textbf{SECS/GEM integration.} Industry-standard fab MES protocol (SEMI E5/E30/E37) via the open-source \texttt{secsgem} Python library (BSD-3, pinned 0.3.0). A Rust adapter service \texttt{mbm-secsgem} ($\sim$\,2000\,LOC) brokers between the binary HSMS wire format and the supervisory REST API. Supports Camstar, Critical Manufacturing, and open-source FabMon MES backends.
  \item \textbf{Operator GUI.} Next.js 15 (App Router, RSC) + React 19 + Tailwind v4 + shadcn/ui + TanStack Query + uPlot (1\,Mpts diagnostic plots) + deck.gl ($10^6$-beam wafer-map heatmap), TypeScript 5.6 strict. Renders on any cleanroom workstation running Firefox ESR 128 or Chromium 129. Page set: \texttt{/} (tool overview, 1\,Hz WebSocket push), \texttt{/lots}, \texttt{/lots/[id]}, \texttt{/recipes}, \texttt{/calibration}, \texttt{/diagnostics}, \texttt{/manual} (gated by hardware YubiKey + physical e-stop interlock), \texttt{/admin}.
  \item \textbf{Alarm system.} Hard alarms (coil quench, vacuum leak, cryocooler trip, ZMI sync loss, He leak) trigger E-stop in firmware; supervisory only records and notifies via Alertmanager $\to$ PagerDuty. Soft alarms (registration drift, dose deviation, source emission drift, tile beam outage) are PromQL queries against VictoriaMetrics on a 1\,s tick, rules version-controlled in \texttt{alarms/*.yaml}, notified via Slack/Mattermost/email.
  \item \textbf{Multi-tenancy + audit.} Keycloak 26 SSO with hardware FIDO2 second factor for process\_engineer and sysadmin roles. HashiCorp Vault 1.18 OSS (or OpenBao 2.0) for secrets. Customer tenancy is row-level-security in Postgres. Audit log is append-only via Postgres trigger; manual-control commands are logged to a write-once tier (S3 Object Lock) for FDA 21 CFR Part 11 forward-compatibility.
  \item \textbf{Reproducible build.} Nix flake pins every tool version (Rust 1.83, Node 22.11, Python 3.13, Postgres 18, VictoriaMetrics 1.110, MinIO RELEASE.2025-10-15). Per-service Docker images via \texttt{dockerTools.buildLayeredImage} (Nix $\to$ OCI), bit-reproducible. K3s for the tool itself, FluxCD 2.4 for GitOps, Sealed Secrets + Vault-CSI for secrets, Harbor 2.12 + Cosign for the container registry.
  \item \textbf{HIL test harness.} \texttt{mbm-hil} simulates the firmware boundary (same Cap'n Proto over UDS), replays canned telemetry, accepts compiled recipe blobs and asserts byte-exact match against fixtures. CI (GitHub Actions self-hosted runners on the build farm) runs the full supervisory stack against \texttt{mbm-hil} on every PR. Wall-time budget: 20 minutes per PR.
\end{itemize}

\textbf{Code skeleton.} Reference repo \texttt{multibeam-ebl-tool-os/} contains a $\sim$50-line Next.js tool-status page (TypeScript), a $\sim$50-line PostgreSQL 18 calibration schema (SQL with UUID v7 and unique-active-snapshot constraint), and a $\sim$50-line Python recipe-to-firmware compiler interface (boto3 + subprocess + content-addressed cache). Full repo layout, role matrix, and 8 open questions for v5.1+ in [S11] \texttt{v5\_software\_operating\_system.md}.

\subsection{Software stack integration}

The three software subsystems exchange three well-defined contracts:

\begin{itemize}[leftmargin=1.6em]
  \item \textbf{Compiler $\to$ firmware}: the binary streaming format in [S9] \S 9.5.5 (\texttt{mbmw\_macrotile} records with BLAKE3 content hashes) is consumed by the per-tile pattern-segment expander ASIC defined in [S10] \S 9.4.
  \item \textbf{Firmware $\to$ OS}: the per-tile RISC-V supervisor (L3) exchanges telemetry and recipe blobs with the warm-side Tango device server (L4) over 10\,GbE TSN via Cap'n Proto, as defined in [S11] \S 0 and [S10] \S 9.7.
  \item \textbf{OS $\to$ compiler}: the recipe-to-firmware compile request (recipe SHA + snapshot UUID + firmware version) is content-addressed cacheable, as defined in [S11] \S 5.2 and consumed by the compiler entry point in [S9] \S 9.11.
\end{itemize}

The three subsystems together close the software loop from customer GDS-II $\to$ blanker plate command at the wafer surface, with every contract defined by a permissively-licensed open-source specification and a runnable reference implementation.

% =====================================================================
\clearpage
\section{Updated quantitative design budgets}

\subsection{Per-beam parameters (v5 nominal)}

\begin{center}
\begin{tabular}{lrr}
\toprule
Parameter & v3 nominal & v5 nominal \\
\midrule
Beam acceleration voltage & 50\,kV & 50\,kV \\
Beam current per channel & 5\,nA & \textbf{2\,nA} \\
Beam spot size at wafer & 5--20\,nm & 5--10\,nm \\
Deflection field strength & 0.1--1\,mT & 0.1--1\,mT \\
Drift distance & 50--150\,mm & 100\,mm \\
Per-beam DAC & 16-bit @ 100\,kHz & 16-bit @ 100\,kHz cryo-CMOS \\
\textbf{Beam pitch (nominal)} & 20\,\textmu m & 20\,\textmu m \\
\textbf{Column operating temperature} & 295\,K (implicit) & \textbf{77\,K (HTS coils)} \\
\textbf{Coil material} & copper microcoil & \textbf{YBCO 2G coated conductor} \\
\textbf{Inner-loop close time (firmware)} & not specified & \textbf{6\,\textmu s of 10\,\textmu s frame} \\
\textbf{Pattern compression cascade} & not specified & \textbf{quadtree + RLE + Zstandard, 10$\times$ logic} \\
\bottomrule
\end{tabular}
\end{center}

\subsection{Aggregate parameters}

\begin{center}
\begin{tabular}{lrr}
\toprule
Parameter & Conservative & Aggressive \\
\midrule
Beam count $N$ & $6 \times 10^4$ & $10^6$ \\
Column array width & 5\,mm & 20\,mm \\
Layer write time (300\,mm wafer, 30\,nm grid, 2\,nA) & 320\,s & 3.2\,s \\
Tool throughput (wph through 70-layer flow) & $\sim$\,5 & $\sim$\,15 \\
Capital cost per tool & \$25\,M & \$42\,M \\
\bottomrule
\end{tabular}
\end{center}

\subsection{Capital cost breakdown (v5 production-grade tool)}

\begin{center}
\begin{tabular}{lrl}
\toprule
Subsystem & Cost share & Source \\
\midrule
Electron source array (MEMS CNT/Spindt) & \$3.0\,M & [S5] \\
Per-beam source closed-loop control electronics & \$0.65\,M & [S5] \\
Acceleration \& condenser optics & \$2.0\,M & (v3 baseline) \\
Per-beam HTS deflection coils + Meissner shields & \$4.7\,M & [S2] \\
Per-beam cryo-CMOS DAC + photonic data path & \$5.3\,M (10-tool amortised) & [S3] \\
Per-beam blanker array (CMOS-MEMS) & \$0.6\,M & [S4] \\
Wafer stage, 3-tier registration loop, ZMI metrology & \$15.0\,M & [S5] \\
Wafer chuck thermal management (AlN + microloop + TEC trim) & \$0.45\,M & [S6] \\
Cryocooler + cold--warm thermal interface (PT-180 + AL-200 dual-zone) & \$1.67\,M & [S2] + [S7] \\
Vacuum + differential pumping + resist mgmt & \$2.3\,M & [S8] \\
Control system + control firmware ASIC IP & \$2.0\,M & (v3 baseline) \\
\textbf{Software dev (compiler + firmware + OS), 10-tool amortised} & \textbf{\$2.5\,M} & \textbf{[S9, S10, S11]} \\
\midrule
\textbf{Total (v5)} & \textbf{$\sim$\,\$40\,M} & \\
\bottomrule
\end{tabular}
\end{center}

Range: \textbf{\$37--42\,M per production tool}. The v5 software NRE adds \$2.5\,M amortised over a 10-tool program. Full undiscounted software NRE is \$2--3\,M one-time across the program: $\sim$\,\$0.8\,M compiler + $\sim$\,\$1.0\,M firmware ASIC software + $\sim$\,\$0.7\,M operating system + $\sim$\,\$0.5\,M integration and HIL. \textbf{Software development is one-time NRE per program, amortised across all tools built}; per-tool the recurring software cost is dominated by deployment and field-engineering support, not development. Versus EUV at \$200--400\,M and DUV at \$50--100\,M, v5 retains the $\sim$\,10$\times$ cost advantage over EUV at comparable nodes.

\subsection{Stochastic placement budget (v5, unchanged from v4.1)}

\begin{center}
\begin{tabular}{lrrr}
\toprule
Contribution & v3 published & v4 calibrated & \textbf{v4.1 / v5 (Niki sim)} \\
\midrule
$\sigma_\text{cross}$ (cross-beam Coulomb) & 5\,nm & 2\,nm at 2\,nA & 2\,nm \\
$\sigma_\text{LJ}$ (Loeffler--Jansen) & 5--20\,nm (band) & 7\,nm & \textbf{6.04\,nm RMS} \\
$\sigma_\text{B}$ (Boersch) & 5\,nm & 3\,nm & \textbf{1.35\,nm} ($C_c = 50$\,\textmu m) \\
$\sigma_\text{stage}$ (registration) & not budgeted & 1\,nm & 1\,nm \\
$\sigma_\text{column}$ (aberration) & not budgeted & 1\,nm & 1\,nm \\
\midrule
\textbf{$\sigma_\text{total}$} & \textbf{17\,nm (cons)} & \textbf{8.0\,nm} & \textbf{6.65\,nm} \\
\bottomrule
\end{tabular}
\end{center}

\subsection{Validated node-range envelope (v5, unchanged from v4.1)}

\begin{center}
\begin{tabular}{rrlll}
\toprule
Node & 20\,\% sub-pixel budget & v3 verdict & v4 verdict & \textbf{v4.1 / v5 verdict} \\
\midrule
180\,nm & 36\,nm & comfortable PASS & comfortable PASS & comfortable PASS \\
90\,nm & 18\,nm & comfortable PASS & comfortable PASS & comfortable PASS \\
\textbf{50\,nm} & \textbf{10\,nm} & \textbf{FAIL} & \textbf{comfortable PASS} & \textbf{comfortable PASS} \\
30\,nm & 6\,nm & FAIL & marginal FAIL & \textbf{MARGINAL} \\
22\,nm & 4.4\,nm & FAIL & FAIL & FAIL \\
\bottomrule
\end{tabular}
\end{center}

\textbf{v5 validated envelope: 30--180\,nm (30\,nm marginal, 50--180\,nm comfortable)}, unchanged from v4.1. The software stack does not alter the placement budget; it implements the budget.

% =====================================================================
\clearpage
\section{Supply-chain check (v5)}

Every subsystem has \textbf{$\geq 3$ independent vendors across $\geq 2$ jurisdictions}, including the v5 software additions. The v4.1 hardware vendor matrix is preserved.

\begin{center}
\begin{tabular}{ll}
\toprule
v5 software subsystem & Library / framework alternatives \\
\midrule
GDS-II / OASIS ingest & gdstk (BSL), gdspy (Boost), klayout (GPL, oracle only) \\
Hierarchical compression & xxHash + zstandard, BLAKE3 + LZMA \\
Pattern-compiler hot loop & Rust via PyO3, Cython, C++ via pybind11 \\
HDL source & Chisel 6 (MIT), Amaranth (BSD) \\
HDL simulation & Verilator (LGPL), cocotb (BSD), Icarus Verilog (GPL) \\
Embedded RTOS & Zephyr (Apache 2), FreeRTOS (MIT) \\
Calibration database & PostgreSQL 18, MinIO, VictoriaMetrics \\
GUI stack & Next.js 15 + React 19 + Tailwind v4 + shadcn/ui \\
MES integration & secsgem (BSD), OpenAS, commercial Camstar/CMF \\
Control framework & Tango Controls (LGPL), EPICS (BSD) \\
Identity / secrets & Keycloak (Apache 2), HashiCorp Vault OSS, OpenBao \\
Container / orchestration & Nix + K3s + FluxCD + Harbor + Cosign \\
\bottomrule
\end{tabular}
\end{center}

\textbf{No new single-vendor chokepoint introduced.} The v5 software stack is built entirely on permissively-licensed open-source libraries with $\geq 2$ maintained alternatives per layer. The only proprietary tools in the chain (Vivado, Cadence Innovus, Synopsys ICC2) are confined to the foundry MPW contract for ASIC P\&R and are isolated behind stable interface contracts (SystemVerilog in, GDS-II / layout out); the firmware repo ships open behavioral models for CI sim so contributors do not need foundry-PDK access.

% =====================================================================
\section{Updated experimental validation pathway}

The v3 staged pathway (Stages 0 $\to$ A $\to$ B $\to$ C $\to$ D) is retained with the v5 software pre-validation inserted as Stage A.5:

\begin{itemize}[leftmargin=1.6em]
  \item \textbf{Stage 0 (numerical)}: Phase 1 cross-beam (complete, v3 \S 3.4); v4.1 Phase 3 engineering sims (complete); Phase 2 GPT/OPAL Loeffler--Jansen (next, highest priority); multi-physics FEA of cryogenic column thermal/magnetic and wafer-chuck transient. \textbf{Estimated cost}: \$0 (open-source codes on commodity Linux + remote compute via cloud GPU as needed).
  \item \textbf{Stage A (single-beam, $\sim$\,\$0.6\,M, 18 months)}: validate Lorentz kinematics with a single HTS coil at 77\,K + cryo-CMOS DAC. Validates \S 3.1 + \S 3.5 simultaneously. Outcome: deflection transfer function, settling time, hysteresis, AC-loss verification.
  \item \textbf{Stage A.5 (software pre-validation, $\sim$\,\$0.3\,M, 3--6 months, new in v5)}: end-to-end software pipeline on simulated firmware / HIL before committing to multi-beam hardware. Compile a representative GDS-II test layout through the v5 pattern compiler (\S 3.14.1); replay the resulting binary stream through the cocotb-emulated firmware on the VCU128 Versal Premium prototype (\S 3.14.2); exercise the operator OS stack against the \texttt{mbm-hil} harness through a full calibration + lot dispatch + abort cycle (\S 3.14.3). \textbf{Stage A.5 gates}: end-to-end latency $<$\,500\,\textmu s tile-to-blanker; compression ratio $\geq$\,10$\times$ on the SPICE-circuit test layout; cocotb regression green; HIL run reproducible across two independent operators. Stage A.5 catches integration bugs before they cost a \$5\,M Stage B respin.
  \item \textbf{Stage B (few-beam prototype, $\sim$\,\$5\,M)}: 4$\times$4 to 8$\times$8 beam array at 20\,\textmu m pitch, full cryogenic column subset, photonic data link (32-channel), cryosorption pumping at proportional density, BSE registration on test fiducials. Validates cross-beam Coulomb, magnetic shielding (Meissner), data path, source closed-loop, and registration. \textbf{The critical de-risking stage.} v5 enhancement: by Stage B the software stack is already field-validated against the Stage A.5 HIL.
  \item \textbf{Stage C (beam-count scaling, \$10--15\,M)}: $10^3$--$10^4$ beams in production-grade column, first wafer exposures at v4 nominal.
  \item \textbf{Stage D (production, \$30--40\,M)}: full $10^5$--$10^6$ beam tool, integration with standard wafer-processing line.
\end{itemize}

Total program cost from Stage 0 through first production article: \textbf{\$45--60\,M, 39--66 month timeline} (Stage A.5 adds 3--6 months at low cost; the resulting reduction in Stage B respin risk typically compresses the integrated schedule rather than extending it).

% =====================================================================
\section{Open engineering questions remaining (v5)}

The v4.1 open-question list (Phase 2 GPT/OPAL sim, YBCO 20\,\textmu m patterning, cryo-CMOS DAC NRE, MEMS CNT/Spindt array scale-up, quench protection, BSE registration crosstalk, Stage B integration) is carried forward. Items that the v5 software stack now closes (or substantially de-risks) are crossed off below; new software-side items are added.

\textbf{Closed (or substantially de-risked) by v5:}

\begin{itemize}[leftmargin=1.6em]
  \item \sout{Pattern-data compiler is the long pole for tool turn-around} (v4 \S 3.6 hand-off note). Closed by [S9]: 3-stage hierarchical compression with runnable code skeleton, FPGA-accelerated ZSTD endpoint, content-addressed binary streaming format.
  \item \sout{Real-time control software architecture not specified} (implicit v4 gap). Closed by [S10]: L0--L4 layered firmware with 6\,\textmu s inner-loop close time, Chisel-generated SystemVerilog, MIT + LGPL open release with Nix-pinned reproducible builds.
  \item \sout{Calibration database, MES integration, operator GUI not specified} (implicit v4 gap). Closed by [S11]: PostgreSQL 18 calibration schema with UUID v7 snapshots and instant-rollback, SECS/GEM adapter via \texttt{secsgem-py}, Next.js 15 operator GUI, Tango Controls supervisor.
\end{itemize}

\textbf{Carried forward from v4.1:}

\begin{enumerate}[leftmargin=1.6em]
  \item \textbf{Phase 2 GPT/OPAL sim} (highest priority, $\sim$\,1 month): rigorous $\sigma_{LJ}$ at v4 geometry with full column optics; expected refinement $\pm 2$\,nm. The v4.1 Jansen-formula evaluation is a first independent check, not a replacement for full particle-in-cell simulation.
  \item \textbf{YBCO 20\,\textmu m patterning}: process development at SuperPower or AMSC, $\sim$\,6--12 months. v4.1 sets a film-quality floor of defect density $\geq 10^{10}$ for the Meissner attenuation requirement.
  \item \textbf{Cryo-CMOS DAC NRE} (Imec + Fraunhofer IPMS dual-source): $\sim$\,12 months.
  \item \textbf{MEMS CNT/Spindt array scale-up} to $10^6$ tips from current $10^4$--$10^5$ state-of-art: $\sim$\,12--18 months.
  \item \textbf{Quench protection at $10^6$-channel scale}: redundancy/detection architecture; $\sim$\,6 months.
  \item \textbf{Per-beam BSE registration crosstalk} at 20\,\textmu m beam pitch: characterisation during Stage B.
\end{enumerate}

\textbf{New in v5 (software-side):}

\begin{enumerate}[leftmargin=1.6em,start=7]
  \item \textbf{Worst-case layer compressibility.} The 10$\times$ cascade target ([S9] \S 9.5.1) is supplier-validated on logic layouts; metrology and test-pattern layers may approach 1$\times$. Open: define a per-layer compressibility metric (Shannon entropy of the rasterised tile distribution) and the threshold that triggers a ``stream-from-DRAM not SSD'' fast path. Joint between [S9] (compiler-side identification) and [S3] / v4 \S 3.6 (transport-side accommodation).
  \item \textbf{Multi-tool farm orchestration.} When a fab runs 4--10 v5 tools ([S11] \S 12.2), the supervisory plane must coordinate load balancing, calibration sharing (where snapshots are tool-portable --- mostly they aren't, coil offsets are per-coil), spare-tool failover mid-lot, and cross-tool yield analytics. Likely add an \texttt{mbm-fleet} service on top of single-tool supervisors with a CRDT-backed state model. Open.
  \item \textbf{RTOS selection bench validation.} Zephyr is selected over FreeRTOS on five rationales ([S10] \S 9.10), but the cryo-MCU port has not been bench-validated. Specific open question: does Zephyr's tickless idle work correctly when the RV32IMC core's clock tree is gated by the cryo-CMOS PLL during low-activity windows? A 1-month bench characterisation on the VCU128 emulator with cryo-trimmed clock margins is the proposed test.
  \item \textbf{Kalman filter divergence under sensor dropout.} If the ZMI interferometer drops out (mechanical bump, fiber breakage), the 100\,Hz Kalman update reverts to BSE-only mode. Open: does the filter remain stable, or does it drift? Standard accelerator-physics answer is ``particle filter with re-initialisation on sensor return''; open whether the warm-side x86 supervisor's compute budget sustains that complexity at 100 Hz across 100 subarrays.
\end{enumerate}

\textbf{No paradigm risk remains.} The architecture's fundamental physics (cross-beam Coulomb, Loeffler--Jansen, cryogenic superconducting deflection, photonic data delivery) is closed within engineering tolerance; the v5 software stack is closed to runnable code skeletons. The remaining work is normal de-risking for any precision-instrument program.

% =====================================================================
\section{Conclusion}

The v3 manuscript introduced a multi-beam direct-write electron lithography architecture based on per-beam Lorentz steering with multi-rate control hierarchy, quantitatively validated the cross-beam Coulomb perturbation envelope, and revised the nominal pitch from 5\,\textmu m to 20\,\textmu m. External peer review of v3 identified seven principal engineering hurdles.

v4 closed each through subsystem engineering (cryogenic column, photonic data path, current de-rating, three-tier registration, AlN chuck thermal, cryosorption). v4.1 confirmed each subsystem with independent Niki-run sims, refining $\sigma_\text{total}$ from 8.0\,nm to 6.65\,nm and exposing the architectural hard floor at 1.41\,nm ($\sim$\,7\,nm node accessibility limit) via the new node-scale curve.

\textbf{v5 closes the software loop.} Three open-source subsystems --- the GDS-II $\to$ blanker-schedule pattern compiler ([S9]), the $10^6$-channel hard real-time control firmware ([S10]), and the tool operating system with calibration database and MES integration ([S11]) --- are specified down to runnable code skeletons, with permissively-licensed library pinning matching the niki-os monorepo conventions. The pattern compiler meets the v4 \S 3.6 10$\times$ compression target via a three-stage cascade (quadtree dedupe + RLE + Zstandard); the firmware closes the inner deflection loop in 6\,\textmu s out of a 10\,\textmu s frame on per-tile cryo-CMOS RISC-V supervisors; the operating system provides a full calibration archive, recipe/lot manager, Next.js operator GUI, SECS/GEM adapter, and Tango-Controls supervisor on a commodity-Linux stack.

Capital cost rises modestly from v4.1's \$37.6\,M to v5's $\sim$\,\$40\,M, with the increment being amortised software NRE across a 10-tool program. Validated node range remains 30--180\,nm (30\,nm marginal). Every subsystem --- hardware and software --- has $\geq 3$ independent vendors or maintained alternatives. The open-supply-chain premise is preserved through both the cryogenic / photonic upgrades and the software stack.

The validation pathway adds Stage A.5 (software pre-validation on HIL, $\sim$\,\$0.3\,M / 3--6 months) between Stage A and Stage B, catching integration bugs at HIL cost rather than at Stage B respin cost. Stage A remains buildable now at \$600\,k / 18 months on commodity laboratory hardware.

The architecture is released openly. The author makes no patent claim on its content. Construction, modification, and commercialisation by any party are encouraged.

\medskip
\noindent\textbf{The architecture is hardware-specified, sim-verified, and software-defined. Stage A is buildable now at \$600\,k / 18 months. The open-source package contains everything needed to begin.}

% =====================================================================
\section*{Acknowledgements}

The author thanks the v3 external reviewer (anonymous, May 2026) whose detailed peer-review report identified the seven engineering hurdles that the v4 manuscript addresses. The v4 subsystem engineering was developed with the assistance of multiple parallel-running AI subagents (Claude general-purpose agents), each tasked with a specific subsystem specification; the human author reviewed, integrated, and is responsible for all engineering decisions and any errors. The Phase 1 numerical simulation of v3 was AI-assisted (Anthropic Claude Code). Phase 3 verification sims (Jansen LJ, cryo thermal network, Meissner attenuation, node-scale curve) were executed by Niki (numpy + matplotlib on Niki box). \textbf{v5 software stack engineered by parallel AI subagents (Claude general-purpose) under human author review. Final integrated open-source release package compiled with tectonic on commodity Linux.}

\section*{Data and Materials Availability}

This is a theoretical proposal with first-pass numerical simulation (Phase 1, v3 supplementary bundle [16]), v4.1 Phase 3 engineering sims, and detailed subsystem engineering specifications (S1--S11). All concepts, equations, figures, architectural choices, engineering specifications, and software reference implementations described herein are released to the public domain via CC BY 4.0 (papers and specs) and MIT (code) licenses. Subsystem specifications (S1--S11) and software reference implementations are released as the open-source design package \texttt{morin\_2026\_multibeam\_ebl\_open}.

\section*{Funding}

This work received no external funding.

\section*{Conflicts of Interest}

The author declares no patent claims on this work and no commercial relationships that would represent conflicts of interest.

% =====================================================================
\section*{Supplementary specifications (companion documents)}

\begin{itemize}[leftmargin=1.6em]
  \item [S0] Morin (2026), ``Multi-Beam EBL v5 Open Source Design Package,'' CC BY 4.0 + MIT. Contains subsystem specifications S1--S11 and software reference implementations.
  \item [S1] \texttt{v4\_loeffler\_jansen.md} --- Loeffler--Jansen calibration from IMS Nano scaling.
  \item [S2] \texttt{v4\_cryocolumn.md} --- Cryogenic column, HTS coils, Meissner shields, cryocooler.
  \item [S3] \texttt{v4\_datapath.md} --- Photonic 400\,Tbps data path, cryo-CMOS DAC ASIC.
  \item [S4] \texttt{v4\_blanker.md} --- Per-beam blanker array (CMOS-MEMS).
  \item [S5] \texttt{v4\_source\_and\_registration.md} --- MEMS field-emitter array + 3-tier registration.
  \item [S6] \texttt{v4\_wafer\_thermal.md} --- Wafer chuck thermal management.
  \item [S7] \texttt{v4\_cold\_warm\_interface.md} --- Cold-warm UHV thermal/magnetic/mechanical interface.
  \item [S8] \texttt{v4\_resist\_outgassing.md} --- Resist selection + outgassing pumping.
  \item [S9] \texttt{v5\_software\_compiler.md} \textbf{(new in v5)} --- GDS-II $\to$ blanker schedule pattern compiler with 3-stage hierarchical compression.
  \item [S10] \texttt{v5\_software\_firmware.md} \textbf{(new in v5)} --- $10^6$-channel hard real-time control firmware (L0--L4 stack on cryo-CMOS).
  \item [S11] \texttt{v5\_software\_operating\_system.md} \textbf{(new in v5)} --- Tool operating system: calibration DB, recipe manager, operator GUI, MES adapter.
\end{itemize}

\section*{References}

References [1]--[23] from v3 and v4.1 are carried forward. Additional v5 software references (tool / library references, not paper citations):

\begin{enumerate}[leftmargin=2.2em,label={[\arabic*]},start=24]
\item L. Heitzmann, ``gdstk: a C++/Python library for creation and manipulation of GDSII and OASIS files,'' BSL-1.0 license, 2020--present. Pattern compiler ingest layer.
\item L. Heitzmann, ``gdspy: Python module for creation and manipulation of GDSII files,'' Boost license, 2011--2021. Legacy ingest layer.
\item M. K\"offerlein, ``KLayout: high-performance EDA viewer and DRC engine,'' GPL-3, 2006--present. DRC oracle (out-of-tree).
\item J. Bachrach et al., ``Chisel: constructing hardware in a Scala embedded language,'' \textit{DAC} 2012; current Chisel 6, MIT license. Parametric HDL generator for the per-tile inner-loop FSM.
\item W. Snyder, ``Verilator: high-performance Verilog simulator,'' LGPL, 1994--present. Cocotb simulation backend for the firmware regression suite.
\item ``cocotb: coroutine-based cosimulation library for writing VHDL and Verilog testbenches in Python,'' BSD, 2013--present. Firmware testbench framework.
\item Zephyr Project, ``Zephyr RTOS,'' Apache 2.0, 2016--present. Per-tile RISC-V supervisor RTOS.
\item PostgreSQL Global Development Group, ``PostgreSQL 18,'' PostgreSQL License, 2025. Calibration database structured metadata store.
\item MinIO Inc., ``MinIO: high-performance, S3-compatible object storage,'' AGPL-v3, 2014--present. Per-beam coefficient blobs + telemetry traces.
\item ``VictoriaMetrics: fast, cost-effective, scalable time series database,'' Apache 2.0, 2019--present. Telemetry time-series ingest.
\item Vercel, ``Next.js 15: the React framework for production,'' MIT, 2016--present. Operator GUI.
\item Tango Controls Collaboration, ``Tango Controls: a control framework for science facilities,'' LGPL, 1999--present. Warm-side L4 supervisor.
\item ``secsgem: Python implementation of SECS/GEM communication protocol,'' BSD-3, 2013--present. SEMI E5/E30/E37 MES adapter.
\item Y. Collet, ``Zstandard: a fast lossless compression algorithm,'' BSD, 2015--present. Entropy stage of the pattern compiler cascade.
\item M. Klein, ``Hierarchical pattern compression for multi-beam mask writers,'' \textit{Proceedings of SPIE} \textbf{11854}, 1185408 (2021). Reference for the IMS-Nano-style hierarchical shot-data format.
\end{enumerate}

\vspace{1em}\hrule\vspace{0.5em}
\begin{center}
\textit{Preprint released openly under CC BY 4.0 (paper + spec) and MIT (code).}\\
\textit{For attribution: Morin, R. G. (2026). ``Multi-Beam Direct-Write Electron Lithography via Multi-Rate Electromagnetic Steering --- v5 Open Software Stack and Final Integrated Release.'' Preprint v5, May 2026.}
\end{center}

\end{document}
